\documentclass[a4paper]{article}
\usepackage{parskip}
\usepackage{lipsum}

\def\nterm {Spring}
\def\nyear {2026}
\def\ncourse {Introduction to Differential Equations}

\input{../header}
\author{Notes taken by \nauthor\vspace{5pt}\\
  Lectures by Gautum Iyer\vspace{5pt}\\
  Carnegie Mellon University}

\newcommand{\TODO}{\textcolor{red}{\textbf{*** TO-DO ***}}}

\begin{document}
\maketitle

\tableofcontents

\section{Introduction to ODEs}

\subsection{Example of ODEs}

An example of ODE is as follows:
\[
  \dot{x} = - \alpha x,
\]
where $\alpha > 0$. This represents exponential decay, whose solution is
  $x(t) = e^{-\alpha t} x_0$.

A more complicated example is the logistic equation where $P(t)$
represents the population density of some species, and we have
\[
  \dot{P} = r P \left( 1 - \frac{P}{K} \right),
\]
where $r, K > 0$. The solution is
\[
  P(t) = \frac{P_0 K e^{rt}}{(K - P_0) + P_0 e^{rt}}.
\]

Another related example is Lotka-Volterra system. Let $x(t)$ be the
population density of the rabbits and $y(t)$ is the population density of
foxes. They can be described by the following model:
\begin{align*}
  \dot{x} & = \alpha x - \beta x y,    \\
  \dot{y} & = - \gamma y + \delta x y,
\end{align*}
where the parameters $\alpha, \beta, \gamma, \delta > 0$.

\subsection{Autonomous linesr systems}

Let $x : [0, \infty) \to \R^d$ and $A \in \R^{d \times d}$. Consider the
equation $\dot{x} = A x$. The solutionon to this equation is $x(t) =
e^{tA} x_0$, where for any square matrix $B$ we defined
\[
  e^B = \sum_{n = 0}^\infty \frac{B^n}{n!}.
\]
It is easy to check that this converges nicely, and thus we can take the
  derivative of $e^{tA}$ against $t$ term by term. We can then verify that
  $x(t) = e^{tA} x_0$ indeed solves the equation.

There are still something to say about this -- how do we compute $e^{tA}$?
If $UAU^{-1} = D$ for some $D$ diagonal, we then have $e^{At} = U^{-1}
e^{Dt} U$. Now if the diagonal matrix $D = \diag(\lambda_1, \dots,
\lambda_d)$, we have $e^{Dt} = \diag(e^{\lambda_1 t, \dots, \lambda_d
t})$.

However, if $A$ is not diagonalizable, we have no option but to consider
the Jordan normal form. Consider a basis change by setting $y = U x$. We
have
\[
  \dot{y} = U \dot{x} = U A x = U A U^{-1} y.
\]
Therefore by choosing matrix $U$ appropriately, we can guarantee
  $UAU^{-1}$ is in Jordan normal form. Hence now we assume $A$ is in
  Jordan normal form. Now decompose $\R^d$ into invariant subspaces
  $\{E_j\}$ such that $\R^d = \bigoplus_j E_j$ and each $E_j$ has
  generalized eigenvalues $\lambda_j = a_j + i b_j$. There are two cases.
\begin{itemize}
  \item If $b_j = 0$, there exists a basis $B_j$ of $E_j$ such that
        \[
          (A)_{B_j} = \begin{bmatrix}
            a_j & 1   &        &     \\
                & a_j & \ddots &     \\
                &     & \ddots & 1   \\
                &     &        & a_j \\
          \end{bmatrix}.
        \]
        Then,
        \[
          (e^{tA})_{B_j} = e^{t a_j} \sum_{k = 0}^{m_j - 1} \frac{t^k}{k!}
          \left( (A)_{B_j} - a_j I \right)^k,
        \]
        where $m_j = \dim(E_j)$.

  \item If $b_j \neq 0$, there exists a basis $B_j$ of $E_j$ such that
        \[
          (A)_{B_j} = \begin{bmatrix}
            C_j & I_2 &        &     \\
                & C_j & \ddots &     \\
                &     & \ddots & I_2 \\
                &     &        & C_j \\
          \end{bmatrix},
        \]
        where $I_2 = \begin{bmatrix}
            1 & 0 \\
            0 & 1
          \end{bmatrix}$
        and $C_j = \begin{bmatrix}
            a_j & -b_j \\
            b_j & a_j
          \end{bmatrix}$.
        Then,
        \[
          (e^{tA})_{B_j} = e^{a_j t} \begin{bmatrix}
            D_j (t) &        &         \\
                    & \ddots &         \\
                    &        & D_j (t)
          \end{bmatrix}  \sum_{k = 0}^{m_j - 1} \frac{t^k}{k!}
          \left( (A)_{B_j} - \diag (C_j, \dots, C_j) \right)^k,
        \]
        where $D_j (t) = \begin{bmatrix}
            \cos (b_j t) & - \sin(b_j t) \\
            \sin (b_j t) & \cos (b_j t)
          \end{bmatrix}$.

\end{itemize}

From the case above, we can see that the long term behvaior of the
solution depends only on the real parts $a_j$. More concretely, let $E^\pm
= \bigoplus_{j \in J^\pm} E_j$ where $J^\pm = \{ j : \sgn(a_j) = \pm 1
\}$, and $E^0 \in \bigoplus_{j \in J^0} E_j$ where $J^0 = \{ j : a_j = 0
\}$. If $x_0 \in E^-$, then $x(t) \to 0$ as $t \to \infty$, and
$\abs{x(t)} \to \infty$ as $t \to -\infty$. This is the \textbf{stable
subspace}. On the other hand, if $x_0 \in E^+$, then $\abs{x(t)} \to
\infty$ as $t \to \infty$, and $x(t) \to 0$ as $t \to -\infty$. This is
the \textbf{unstable subspace}.

\subsection{Time inhomogenous linear equations}

Consider $\dot{x} = A(t) x$ where $A : \R \to \R^{d \times d}$. In the 1
dimensional case, we have $\dot{x} = a(t) x$. With $x(t_0) = x_0$, the
solution to this equation is
\[
  x(t) = \exp \left( \int_{t_0}^t a(s) \d s \right) x_0.
\]

\begin{thm}
  Fix $t_0, T$. If for each $t_1, t_2 \in [t_0, T]$ we have
  $[A(t_1), A(t_2)] = 0$, then
  \[
    x(t) = \exp \left( \int_{t_0}^t A(s) \d s \right) x_0
  \]
  solves $\dot{x} = A(t) x$ with $x(t_0) = x_0$.
\end{thm}

\begin{proof}
  Let $B(t) = \int_{t_0}^{t} A(s) \d s$. Compute the derivative with
  the Wilcox Formula
  \[
    \dd_t e^{B(t)} = \int_0^1 e^{\alpha B(t)} A(t) e^{(1 - \alpha) B(t)}
    \d \alpha. \tag{W}
  \]
  Note that (W) implies the theorem, because assumption implies that
  $[A(t), B(t)] = 0$ for all $t$. It then follows that $[A(t), e^{\alpha
  B(t)}] = 0$ and thus
  \[
    \dd_t e^{B(t)} = \int_0^1 e^{B(t)} A(t) \d \alpha = A(t) e^{B(t)}.
  \]

  As for the proof for (W), the trick is to define $g(\alpha, t) = e^{-
  \alpha B(t)} \dd_t e^{\alpha B(t)}$. Then we compute
  \begin{align*}
    \dd_\alpha g
     & = - B e^{-\alpha B} \dd_t e^{\alpha B}
    + e^{-\alpha B} \dd_t (B e^{\alpha B})                                   \\
     & = - B e^{-\alpha B} \dd_t e^{\alpha B} + e^{-\alpha B} A e^{\alpha B}
    + e^{-\alpha B} B \dd_t e^{\alpha B}                                     \\
     & = e^{-\alpha B} A e^{\alpha B}.
  \end{align*}
  Then, $g(1) - g(0) = \int_0^1 \dd_\alpha g \d \alpha$. Therefore,
  $e^{-B} \dd_t e^{B} = \int_0^1 e^{-\alpha B} A e^{\alpha B} \d \alpha$,
  and thus
  \[
    \dd_t e^B = \int_0^1 e^{(1 - \alpha) B} A e^{\alpha B},
  \]
  as desired.

\end{proof}

Note that the equation $\dot{x} = A(t) x$ is linear, so any linear
combination of solutions is again a solution. Say $\phi_1, \dots, \phi_d$
are $d$ solutions, and $\phi_i (t_0) = e_i$. Then, given $x \in \R^d$,
$x(t) = \sumi^d x_i \phi_i (t)$, where $x_i = x_0 \cdot e_i$, solves the
equation with initial condition $x_0$. Now define
\[
  \Pi (t_0, t) = \begin{bmatrix}
    \phi_1 (t) & \cdots & \phi_d (t)
  \end{bmatrix}
\]
as the principle matrix. Note that $\dd_t \Pi (t_0, t) = A(t) \Pi (t_0,
  t)$, and $\Pi (t_0, t_0) = I$. Now the solution to $\dot{x} = A(t) x$
  with $x(t_0) = x_0$ is given by $x(t) = \Pi(t, t_0) x_0$.

\begin{prop}
  $\Pi (s, t) \Pi (r, s) = \Pi (r, t)$.
\end{prop}
\begin{proof}
  Fix $r, s$ and let $L(t) = \LHS$ and $R(t) = \RHS$. Note
  that $\dd_t L = A(t) \Pi(s, t) \Pi(r, s) = A(t) L(t)$. Similarly, $\dd_t
    R = A(t) R(t)$. For an initial data, note that $L(s) = R(s) = \Pi (r,
    s)$.
\end{proof}

\begin{cor}
  $\Pi(s, t) = \Pi(t, s)^{-1}$.
\end{cor}

\begin{prop}[Liouville formula]
  $\dd_t (\det (\Pi (s, t))) = \tr(A) \det(\Pi (s, t))$.
\end{prop}

\begin{proof}
  Exercise.
\end{proof}

\section{Existence and Uniqueness}

Now we want to discuss existence and uniqueness of solution for $\dot{x} =
f(t, x)$.
\begin{thm}
  Let $U \subset \R^{d + 1}$ open, $f \in C(U, \R^d)$. Assume $f$
  locally Lipshitz. Then, given $(t_0, x_0) \in U$, there exists a locally
  unique solution to $\dot{x} = f(t, x)$ with $x(t_0) = x_0$.
\end{thm}
By locally unique, we mean a solution in an interval around $t_0$.

\end{document}